\begin{definition}
A \emph{complex-gradient package} is a datum consisting of proofs of the
following two statements.

\emph{First statement.}  Let $\Omega\subset\mathbb{C}$ be open, let $1<p<\infty$, let
$u\colon\mathbb{C}\to\mathbb{R}$ be $p$-harmonic on $\Omega$ in the sense of
\noderef{PHarmonicOn}, and let $g\colon\mathbb{C}\to\mathbb{C}$ be such that
$\operatorname{Re}g$ and $\operatorname{Im}g$ are the distributional
$\partial/\partial x$- and $\partial/\partial y$-derivatives of $u$ on
$\Omega$, in the sense of \noderef{IsWeakPartialOn}; that is, $g=u_x+iu_y$
records the weak gradient $\nabla u$.  Then the complex gradient
\[
  f=\frac{\partial u}{\partial z}=\tfrac12\bigl(u_x-iu_y\bigr)=\tfrac12\,\overline{g}
\]
satisfies
\begin{enumerate}
\item[(a)] $f\in W^{1,2}_{\mathrm{loc}}(\Omega)$ in the sense of
\noderef{MemWLocOn};
\item[(b)] $f$ satisfies equation (12) on $\Omega$ in the sense of
\noderef{SolvesEq12}; and
\item[(c)] $f$ is $K$-quasiregular on $\Omega$ in the sense of
\noderef{QuasiregularOn}, with
\[
  K=\max\Bigl\{p-1,\frac{1}{p-1}\Bigr\},
\]
the exponent of equation (13).
\end{enumerate}

\emph{Second statement.}  Let $\Omega$, $p$, $u$ and $g$ be as above, and let
$f_c\colon\mathbb{C}\to\mathbb{C}$ satisfy $f_c=f$ almost everywhere on
$\Omega$ and be continuous on $\Omega$.  Then
\begin{enumerate}
\item[(d)] $f_c$ is real-differentiable at almost every point of $\Omega$; and
\item[(e)] $f_c$ satisfies equation (12) on $\Omega$ in the sense of
\noderef{SolvesEq12}.
\end{enumerate}

This is item (C1) of the paper's note for the formalization, which quotes it
from [BI] and [Ma]: ``If $u\in W^{1,p}_{\mathrm{loc}}(\Omega)$,
$1<p<\infty$, is $p$-harmonic on a plane domain $\Omega$, then its complex
gradient $f=\partial u/\partial z$ belongs to $W^{1,2}_{\mathrm{loc}}(\Omega)$
and satisfies (12) almost everywhere; in particular $f$ is $K$-quasiregular
with $K$ as in (13).''  The note directs that such quoted results be carried
as explicit hypotheses, bundled into a structure or passed as arguments, and
that they not be introduced as axioms; this definition is that bundle, so that
a theorem needing (C1) takes a complex-gradient package as an argument.  The
hypothesis $u\in W^{1,p}_{\mathrm{loc}}(\Omega)$ of the quoted statement is
part of $p$-harmonicity as \noderef{PHarmonicOn} defines it.

The second statement reads the same quoted sentence for a continuous
representative of the complex gradient.  It is not a strengthening of the
first in its clause (e): equation (12) as \noderef{SolvesEq12} states it is an
identity between the \emph{weak} Wirtinger derivatives of its argument, and by
\noderef{IsWeakWirtingerOn} those depend only on the almost-everywhere class
of that argument on $\Omega$ --- the local integrability and the
integration-by-parts clauses see the argument only through integrals, and the
pointwise values of $f_c$ occur in (12) only inside an ``almost every $z$''
quantifier.  So (e) for $f_c$ and (b) for $\tfrac12\bar g$ are equivalent,
given $f_c=f$ almost everywhere, and (e) is recorded only to spare the
consumer that unfolding.  This is the right behaviour: the paper's complex
gradient is the continuous object --- paper.tex:156 states that
$f=\partial u/\partial z$ is continuous and writes its zero set as
$f^{-1}(0)=\{z\in\Omega:\nabla u(z)=0\}$ --- and the quoted assertion that $f$
``satisfies (12) almost everywhere'' should hold of that representative as
well as of the one attached to a chosen weak gradient $g$, which with the weak
reading of (12) it does.

The one item of the second statement that the quoted sentence does not contain
is clause (d), the almost-everywhere real differentiability of the continuous
representative.  It is quoted from the same sources and in the same spirit as
the rest: a quasiregular mapping is differentiable almost everywhere, by the
$W^{1,2}_{\mathrm{loc}}$ membership of clause (a) together with the standard
regularity theory of [Ah], [LV], and clause (c) makes $f$ quasiregular.  It is
carried here rather than proved because it is an \emph{input} to the
elliptic-regularity bundle (C3), which assumes it of the representative whose
smoothness it concludes, so it cannot be obtained from (C3); and it is not
among the three conclusions that (C2) quotes for the inverse of a
quasiconformal homeomorphism.  The representative itself is not produced here:
the consumer supplies it, obtaining it from the continuity clause of (C2).
\end{definition}
